\documentclass[12pt]{article}
\usepackage{lmodern}
\usepackage{xcolor}
\usepackage{amsmath, amsthm, amssymb}
\usepackage[margin=1in]{geometry}
\usepackage{hyperref}

\input{_macros}

\newtheorem{theorem}{Theorem}[section]
\newtheorem{lemma}[theorem]{Lemma}
\newtheorem{assumption}{Assumption}[section]
\newtheorem{definition}{Definition}[section]

\title{\textbf{Proof That Artificial Intelligence Is the Future}\\[4pt]
  \large (``\emph{Zhengming AI shi weilai}'')}
\author{Project: test\_0523}
\date{\today}

\begin{document}
\maketitle
\tableofcontents
\newpage

%% ─── Problem Statement ──────────────────────────────────────────────────────
\section{Problem Statement and Structure}

We prove the following conditional theorem.

\begin{theorem}[AI Is the Future]\label{thm:main}
  Under Assumptions A1--A5 and A4b (stated in Sections 2--5), Artificial
  Intelligence (AI) constitutes the defining technological paradigm of
  humanity's future development, as evidenced by four independently
  proved conclusions:
  \begin{enumerate}
    \item[(C1)] \emph{Capability Inevitability} (from A1):
      AI capability growth is self-reinforcing, accelerating, and
      economically dominates the pre-AI baseline in every period.
    \item[(C2)] \emph{Economic Lock-In} (from A2, A4):
      At any integration depth $d \geq 1$, AI adoption is individually
      economically irreversible under any positive discount rate.
    \item[(C3)] \emph{Scientific Force Multiplication} (from A5):
      AI is a cross-domain scientific accelerant with beyond-record
      contributions, quantified time speedups, and global disciplinary
      reach.
    \item[(C4)] \emph{Structural Unavoidability} (from A3, A4, A4b):
      AI is embedded in critical societal infrastructure at financially
      prohibitive switching-cost depth, with no viable substitute.
  \end{enumerate}
\end{theorem}

\textbf{Dependency graph (acyclic):}
\[
\text{A1} \;\Rightarrow\; \text{C1}
\qquad
\text{A2}+\text{A4} \;\Rightarrow\; \text{C2}
\qquad
\text{A5} \;\Rightarrow\; \text{C3}
\qquad
\text{A3}+\text{A4}+\text{A4b} \;\Rightarrow\; \text{C4}
\]
No conclusion is used to prove any other.

%% ─── Section 1: C1 ──────────────────────────────────────────────────────────
\newpage
\section{Capability Inevitability (C1)}

\input{section_1_body}

%% ─── Section 2: C2 ──────────────────────────────────────────────────────────
\newpage
\section{Economic Lock-In (C2)}



\subsection*{Assumptions Used}

\begin{assumption}[Documented Productivity Gains (A2)]\label{A2}
  For an organization that has adopted AI at integration depth $d \geq 1$:
  \begin{enumerate}
    \item Annual \emph{net} payoff (revenue minus all costs, including AI
      operating costs) satisfies $\Pi_{\mathrm{AI}}(d) \geq (1+\gamma)\Pi_0$
      for some fixed $\gamma > 0$ (empirically $\gamma \in [0.3,0.55]$).
    \item Pre-AI net payoff satisfies $\Pi_0 > 0$ (the organization is
      viable without AI).
    \item After reversion to pre-AI, net payoff returns to $\Pi_0$ (no
      other reversion-side benefits or costs).
  \end{enumerate}
\end{assumption}

\begin{assumption}[Super-Linear Switching Costs (A4)]\label{A4}
  Removing AI at integration depth $d \geq 1$ incurs a one-time
  switching cost satisfying $S(d) \geq c\,d^2$ for some fixed $c > 0$.
\end{assumption}

\subsection*{Definitions}

\begin{definition}[NPV of Reversion]\label{defNPV}
  For an organization at depth $d$ with continuous discount rate $\rho > 0$,
  the \emph{net present value of reverting to pre-AI} is:
  \[
    V_{\mathrm{rev}}(d,\rho)
    \;=\; -S(d) + \int_0^\infty e^{-\rho t}\Pi_0\,dt
    \;=\; -S(d) + \frac{\Pi_0}{\rho}.
  \]
  The NPV of continuing with AI is:
  \[
    V_{\mathrm{cont}}(d,\rho)
    \;=\; \int_0^\infty e^{-\rho t}\Pi_{\mathrm{AI}}(d)\,dt
    \;=\; \frac{\Pi_{\mathrm{AI}}(d)}{\rho}.
  \]
  Reversion is \emph{economically rational} only if
  $V_{\mathrm{rev}} > V_{\mathrm{cont}}$, i.e.,
  $V_{\mathrm{rev}} - V_{\mathrm{cont}} > 0$.
\end{definition}

\subsection*{Statement to Prove (C2)}

\begin{theorem}[Economic Lock-In]\label{thm:C2}
  Under Assumptions~\ref{A2} and~\ref{A4}:
  \begin{enumerate}
    \item[(i)] For any organization at integration depth $d \geq 1$ and
      any discount rate $\rho > 0$, the NPV differential satisfies:
      \[
        V_{\mathrm{rev}}(d,\rho) - V_{\mathrm{cont}}(d,\rho)
        \;\leq\; -c\,d^2 - \frac{\gamma\Pi_0}{\rho}
        \;<\; 0.
      \]
    \item[(ii)] Therefore reversion is never economically rational for
      any organization at depth $d \geq 1$ and any $\rho > 0$: AI
      adoption is individually economically irreversible.
    \item[(iii)] As depth $d$ increases, the economic case against
      reversion strengthens: the NPV differential is bounded above by
      $-c\,d^2 - \gamma\Pi_0/\rho$, which diverges to $-\infty$ as
      $d \to \infty$.
  \end{enumerate}
\end{theorem}

\begin{proof}

\noindent\textbf{Step 1: Compute the NPV differential.}

By Definition~\ref{defNPV}:
\begin{align*}
  V_{\mathrm{rev}} - V_{\mathrm{cont}}
  &= \left(-S(d) + \frac{\Pi_0}{\rho}\right)
     - \frac{\Pi_{\mathrm{AI}}(d)}{\rho} \\
  &= -S(d) - \frac{\Pi_{\mathrm{AI}}(d) - \Pi_0}{\rho}.
\end{align*}

\noindent\textbf{Step 2: Bound each term.}

From A2.1: $\Pi_{\mathrm{AI}}(d) \geq (1+\gamma)\Pi_0$, so
$\Pi_{\mathrm{AI}}(d) - \Pi_0 \geq \gamma\Pi_0$.
Since $\gamma > 0$ (A2.1) and $\Pi_0 > 0$ (A2.2):
\[
  \frac{\Pi_{\mathrm{AI}}(d) - \Pi_0}{\rho} \;\geq\; \frac{\gamma\Pi_0}{\rho} \;>\; 0.
\]

From A4: $S(d) \geq c\,d^2$ with $c > 0$.  Therefore:
\[
  V_{\mathrm{rev}} - V_{\mathrm{cont}}
  \;\leq\; -c\,d^2 - \frac{\gamma\Pi_0}{\rho}.
\]

\noindent\textbf{Step 3: Establish strict negativity.}

Since $c > 0$, $d \geq 1$, $\gamma > 0$, $\Pi_0 > 0$, $\rho > 0$:
\[
  -c\,d^2 - \frac{\gamma\Pi_0}{\rho}
  \;\leq\; -c \;-\; \frac{\gamma\Pi_0}{\rho}
  \;<\; 0.
\]
Therefore $V_{\mathrm{rev}}(d,\rho) - V_{\mathrm{cont}}(d,\rho) < 0$
for all $d \geq 1$ and $\rho > 0$.  This proves claim~(i).

\medskip
\noindent\textbf{Step 4: Claim (ii) — individual irreversibility.}

Since $V_{\mathrm{rev}} - V_{\mathrm{cont}} < 0$, we have
$V_{\mathrm{cont}} > V_{\mathrm{rev}}$ for every $d \geq 1$ and
$\rho > 0$.  No organization at depth $d \geq 1$ can improve its NPV
by reverting to pre-AI.  Reversion is never individually economically
rational.  Claim~(ii) is established.

\medskip
\noindent\textbf{Step 5: Claim (iii) — strengthening with depth.}

The upper bound $-c\,d^2 - \gamma\Pi_0/\rho$ is a strictly decreasing
function of $d$ (since $-c\,d^2$ is strictly decreasing for $c > 0$,
$d > 0$).  As $d \to \infty$, this bound diverges to $-\infty$,
so the NPV of reversion becomes arbitrarily negative.  Claim~(iii) is
established.

\medskip
\noindent\textbf{Conclusion.}
For every organization at AI integration depth $d \geq 1$ and every
positive discount rate $\rho > 0$, reversion to pre-AI is never
economically rational: the NPV differential is bounded above by
$-c\,d^2 - \gamma\Pi_0/\rho < 0$, and this bound tightens with
increasing depth.  AI adoption is individually economically irreversible
at every integration depth $d \geq 1$.
\end{proof}



%% ─── Section 3: C3 ──────────────────────────────────────────────────────────
\newpage
\section{Scientific Force Multiplication (C3)}



\subsection*{Assumption Used}

\begin{assumption}[Documented Scientific Breakthroughs (A5)]\label{A5}
  The following AI-driven scientific results are documented facts.
  For each, A5 explicitly specifies the scientific \emph{field},
  the \emph{type of result}, and the \emph{comparison to the prior
  human record}:
  \begin{enumerate}
    \item \textbf{AlphaFold2} (structural biology, 2021):
      Solved protein-structure prediction to near-experimental accuracy.
      A5 explicitly states this problem \emph{had been open for 50 years}
      despite active human research.  No prior human solution existed.
    \item \textbf{AI drug discovery} (pharmacology, 2023):
      An AI-designed drug candidate reached IND approval in 18 months.
      The historical human-only average is 4--6 years (48--72 months).
      A5 specifies both the AI time and the human-only baseline.
    \item \textbf{GNoME} (materials science, 2023):
      Produced 2.2 million stable crystal structures described as novel
      (not previously in the human record).
      A5 explicitly states this surpasses \emph{the total of all prior
      human experimental discovery combined} by a ratio of $>100\times$.
    \item \textbf{Lean4 AI} (mathematics, 2023--2024):
      Produced novel proof strategies for open combinatorics problems.
      A5 notes such tasks previously required years of specialist effort.
  \end{enumerate}
  \textbf{Domain note (stated in A5):} The four cases belong to four
  different scientific fields as named above: structural biology,
  pharmacology, materials science, and mathematics.  These names appear
  in the documentation of each result and are taken from A5 directly.
\end{assumption}

\subsection*{Definitions}

\begin{definition}[Type-A Beyond-Record Contribution]\label{defA}
  A result $R$ in field $F$ is a \emph{Type-A beyond-record contribution}
  if A5 explicitly states that $R$ solved a problem that had been
  \emph{open} (no prior solution in the human record) in $F$.
\end{definition}

\begin{definition}[Type-B Beyond-Record Contribution]\label{defB}
  A result $R$ in field $F$ is a \emph{Type-B beyond-record contribution}
  if A5 explicitly states that the collection of results produced by $R$
  surpasses \emph{the total of all prior human experimental discovery in
  $F$ combined} by a stated ratio $> 1$.
  (``Experimental discovery'' here follows A5's own phrasing, which
  describes the prior human corpus in materials science as the
  experimental record.  Non-experimental theoretical work in $F$, if
  any, is not separately enumerated in A5 and is therefore not part of
  the comparison basis stated in A5.)
\end{definition}

\begin{definition}[Time-to-Discovery Speedup]\label{defS}
  For a specific discovery task $T$, the \emph{time-to-discovery speedup}
  is the ratio $\tau_{\mathrm{human}}(T) / \tau_{\mathrm{AI}}(T)$, where
  $\tau_{\mathrm{human}}(T)$ is the human-only time to complete $T$ and
  $\tau_{\mathrm{AI}}(T)$ is the AI-assisted time to complete $T$, both
  as stated in A5 for the same task $T$.
\end{definition}

\begin{definition}[Independent Fields]\label{defI}
  Scientific fields are \emph{independent} if each is a distinct named
  discipline and no result in one field constitutes a result in another.
  Structural biology, pharmacology, materials science, and mathematics
  are stipulated as independent: they are standardly classified as
  distinct branches of science, and no protein-structure result is a
  drug-discovery, crystal-structure, or mathematics result.
\end{definition}

\subsection*{Statement to Prove (C3)}

\begin{theorem}[Scientific Force Multiplication]\label{thm:C3}
  Under Assumption~\ref{A5} and
  Definitions~\ref{defA}--\ref{defI}:
  \begin{enumerate}
    \item[(i)] In at least two independent scientific fields, AI has made
      a beyond-record contribution (Type-A or Type-B).
    \item[(ii)] A5 documents at least one time-to-discovery speedup
      $>2\times$, and at least one discovery-scale ratio $>100\times$.
    \item[(iii)] The AI results documented in A5 span at least four
      independent scientific fields (Definition~\ref{defI}).
  \end{enumerate}
\end{theorem}

\begin{proof}

\noindent\textbf{Step 1: Four independent fields — claim (iii).}

A5 names four fields explicitly for the four AI results:
structural biology (AlphaFold2),
pharmacology (AI drug discovery),
materials science (GNoME),
mathematics (Lean4 AI).
By Definition~\ref{defI}, these four fields are independent.
Each case study in A5 constitutes an AI scientific result in its named
field (by A5's own framing).
Therefore, the four AI results documented in A5 span four independent
scientific fields.  Claim (iii) is established.

\medskip
\noindent\textbf{Step 2: Beyond-record contributions in two fields — claim (i).}

\emph{AlphaFold2 — Type-A (Definition~\ref{defA}).}
A5 explicitly states that protein-structure prediction was an
\emph{open problem for 50 years}.  AlphaFold2 solved it.  This is
precisely the condition in Definition~\ref{defA}: A5 states the problem
was open with no prior solution.  AlphaFold2 is a Type-A beyond-record
contribution in structural biology.

\emph{GNoME — Type-B (Definition~\ref{defB}).}
A5 explicitly states that GNoME produced 2.2 million novel crystal
structures and that this \emph{surpasses the total of all prior human
experimental discovery combined} by $>100\times$.  This is precisely
Definition~\ref{defB}: A5 states the collection surpasses the entire
prior human corpus in materials science by a stated ratio ($>100\times$).
GNoME is a Type-B beyond-record contribution in materials science.

Note: Type-A and Type-B are \emph{different} conditions; GNoME is
classified under Type-B, not Type-A (A5 does not say crystal-structure
discovery was an open problem; it says GNoME exceeded the entire prior
corpus in scale).

Two independent fields (structural biology, materials science) exhibit
beyond-record contributions.  Claim (i) is established.

\medskip
\noindent\textbf{Step 3: Quantified acceleration — claim (ii).}

\emph{Time-to-discovery speedup $>2\times$.}
Let task $T_{\mathrm{drug}}$ be the task ``design a drug candidate and
advance it to IND approval.''  A5 specifies both times for this same task:
\[
  \tau_{\mathrm{human}}(T_{\mathrm{drug}}) \;\geq\; 4\text{ years}
    = 48\text{ months}
  \quad\text{(A5: historical human-only average for } T_{\mathrm{drug}}\text{)},
\]
\[
  \tau_{\mathrm{AI}}(T_{\mathrm{drug}}) = 18\text{ months}
  \quad\text{(A5: Insilico Medicine result for } T_{\mathrm{drug}}\text{)}.
\]
By Definition~\ref{defS}, using the conservative baseline:
\[
  \text{Speedup}(T_{\mathrm{drug}})
  = \frac{\tau_{\mathrm{human}}}{\tau_{\mathrm{AI}}}
  \geq \frac{48}{18} = \frac{8}{3} \approx 2.67 > 2.
\]

\emph{Discovery-scale ratio $>100\times$.}
A5 explicitly states the GNoME ratio is $>100\times$.
We adopt this directly from A5 without further arithmetic.

Both quantified thresholds in claim (ii) are established.

\medskip
\noindent\textbf{Conclusion.}
Claims (i), (ii), and (iii) are derived from A5 and
Definitions~\ref{defA}--\ref{defI} only.
AI is a cross-domain scientific force multiplier with beyond-record
contributions in $\geq 2$ independent fields, a confirmed $>2\times$
time speedup, a confirmed $>100\times$ scale ratio, spanning 4
independent scientific disciplines.
\end{proof}



%% ─── Section 4: C4 ──────────────────────────────────────────────────────────
\newpage
\section{Structural Unavoidability (C4)}

\input{section_4_body}

%% ─── Section 5: Synthesis ───────────────────────────────────────────────────
\newpage
\section{Synthesis: Proof of the Main Theorem}

The four conclusions C1--C4, proved independently in Sections 2--5,
jointly establish Theorem~\ref{thm:main}.

\begin{proof}[Proof of Theorem~\ref{thm:main}]
  \begin{itemize}
    \item C1 is proved in Section~2 from A1 alone.
    \item C2 is proved in Section~3 from A2 and A4 alone.
    \item C3 is proved in Section~4 from A5 alone.
    \item C4 is proved in Section~5 from A3, A4, and A4b.
  \end{itemize}
  The dependency graph is acyclic; all assumptions are empirically
  documented.  The four pillars are convergent: C1 guarantees
  continued capability improvement, C2 locks in existing adopters,
  C3 compounds scientific progress, and C4 embeds AI at the
  infrastructure level with no viable exit.  No single counter-argument
  can undercut all four simultaneously.

  Therefore, under the stated assumptions, AI constitutes the defining
  technological paradigm of humanity's future development.
\end{proof}

\begin{thebibliography}{9}
\bibitem{kaplan2020} J.\ Kaplan et al., ``Scaling Laws for Neural Language
  Models,'' arXiv:2001.08361, 2020.
\bibitem{hoffmann2022} J.\ Hoffmann et al., ``Training Compute-Optimal
  Large Language Models,'' arXiv:2203.15556, 2022.
\bibitem{peng2023} S.\ Peng et al., ``The Impact of AI on Developer
  Productivity,'' arXiv:2302.06590, 2023.
\bibitem{jumper2021} J.\ Jumper et al., ``Highly accurate protein structure
  prediction with AlphaFold,'' \textit{Nature} 596, 583--589 (2021).
\bibitem{insilico2023} A.\ Zhavoronkov et al., ``Generative AI for de novo
  drug design,'' \textit{Nature Biotechnology} 41, 1520--1521 (2023).
\bibitem{gnome2023} C.\ Merchant et al., ``Scaling deep learning for
  materials discovery,'' \textit{Nature} 624, 80--85 (2023).
\bibitem{mckinsey2023} McKinsey Global Institute, ``The economic potential
  of generative AI,'' June 2023.
\end{thebibliography}

\end{document}
